% =============================================================================
% cap06.tex — MODELAGEM E ANALISE ESTRUTURAL
% =============================================================================

\chapter{MODELAGEM E ANÁLISE ESTRUTURAL}\label{cap:modelagem}

\section{SISTEMA ESTRUTURAL}\label{sec:sistema-estrutural}

Cada pórtico transversal é constituído por duas vigas de duplo-bisel apoiadas
sobre três pilares. As vigas são monolíticas, contínuas na seção de cumeeira, e
recebem apoio articulado nas extremidades. Os pilares são articulados na base e
no topo, de modo que não transmitem momento à fundação nem à viga.

A estabilidade lateral não é, portanto, obtida por continuidade de nós, e sim
pelos sistemas de contraventamento: o diafragma de cobertura conduz as forças
transversais aos frontões, e os painéis de cruz nas fachadas longitudinais
conduzem as forças longitudinais à fundação. Essa separação entre função
gravitacional e função de estabilização é a que permite manter os pilares
articulados e com seção reduzida.

\subsection{Modelo do apoio horizontal}\label{ssec:modelo-apoio}

A definição mais consequente da modelagem diz respeito ao apoio horizontal da
viga. Uma viga de duplo-bisel apoiada nas extremidades tem comportamento
intermediário entre dois extremos. Se as extremidades forem impedidas de afastar
horizontalmente, a peça trabalha como arco abatido: desenvolve empuxo e o
momento fletor é aliviado pelo efeito de arco. Se puderem afastar livremente, a
peça trabalha como viga biapoiada, sem empuxo, e o momento é integralmente
resistido por flexão.

A distinção não é acadêmica. Entre um extremo e outro o momento na cumeeira
varia em mais de setenta por cento, e o esforço horizontal transmitido à
fundação varia de zero a algumas centenas de quilonewtons. Adotar o modelo
errado conduz ou a subdimensionar a viga, ou a dimensionar uma fundação para um
esforço que não existe.

O comportamento efetivo é determinado pela rigidez relativa entre o mecanismo
que gera o empuxo e o que o restringe. A viga, como arco, opõe-se ao afastamento
com rigidez de 208.462 kN/m. O pilar que a apoia opõe-se com rigidez nula, por
ser articulado na base, e ainda que fosse engastado ofereceria apenas
52,794 kN/m, valor 3.949 vezes menor que o da viga. O empuxo só se desenvolve na
medida em que a restrição é comparável à rigidez do arco, e não é o caso.

A \autoref{tab:rigidez-apoio} apresenta a varredura da rigidez do apoio, do
engastamento perfeito à liberdade total, obtida pelo método das forças conforme
o item 5.4 da EN 1995-1-1 \cite{EN1995-1-1:2004} e o item 5.2.1 da
EN 1993-1-1 \cite{EN1993-1-1:2005}. Os valores devem ser lidos contra os
208.462 kN/m da rigidez do arco.

\tabela{Empuxo desenvolvido e momento na viga em função da rigidez do apoio}{\footnotesize
\begin{tabular}{rrrl}
\toprule
    \textbf{Rigidez do apoio} & \textbf{Empuxo} & \textbf{Momento na} & \textbf{Situação} \\
    \textbf{(kN/m)} & \textbf{(kN)} & \textbf{viga (kN$\cdot$m)} & \\
\midrule
    Infinita      & 266,85 & 294,2 & Arco perfeito \\
    1.000.000     & 220,82 & 326,4 & \\
    100.000       &  86,51 & 430,0 & \\
    52,794        &   0,07 & 504,2 & Pilar engastado na base \\
\midrule
    \textbf{0}    &   \textbf{0,00} & \textbf{504,3} & \textbf{Pilar articulado, adotado} \\
\bottomrule
\end{tabular}}{Elaborado pelo autor (2026)}{rigidez-apoio}{}{}

A leitura da tabela é inequívoca. Mesmo no limite generoso da base engastada,
desenvolve-se 0,0\% do empuxo rígido, isto é, 0,07 kN de um total possível de
266,85 kN. A restrição oferecida pela base não alcança o apoio da viga, porque
o pilar se deforma por flexão em série com ela.

Adotou-se, em consequência, o modelo de viga biapoiada, com empuxo nulo. O
momento de projeto na cumeeira passa de 294,2 kN$\cdot$m, que o modelo de arco
indicaria, para 504,3 kN$\cdot$m, um acréscimo de 71,4\%. Trata-se do modelo
conservador quanto à flexão, e é também o fisicamente correto.

Duas consequências de projeto decorrem dessa decisão. A primeira é que não há
esforço horizontal a transmitir à fundação por efeito de arco, o que dispensa
tirante de base e simplifica a sapata. A segunda é que a hipótese, aventada em
fase anterior do trabalho, de que o empuxo produziria momento da ordem de
1.300 kN$\cdot$m nos pilares e levaria ao seu colapso, fica afastada: o pilar
não atrai o empuxo.

\section{AÇÕES}\label{sec:acoes}

\subsection{Ações permanentes}\label{ssec:acoes-permanentes}

O peso próprio não foi arbitrado por valor único de área. Cada parcela é
calculada a partir da geometria e da massa específica média do material, de modo
que uma alteração de seção se propague automaticamente ao carregamento. A
\autoref{tab:peso-proprio} apresenta a decomposição, conforme o item 4.1 da
EN 1991-1-1 \cite{EN1991-1-1:2002}, com massa específica da
EN 14080 \cite{EN14080:2013}. A área tributária de um pórtico intermediário é de
5,75 m, igual ao espaçamento entre pórticos, e as parcelas declaradas constam da
\autoref{tab:valores-declarados}.

\tabela{Decomposição da ação permanente característica}{\footnotesize
\begin{tabular}{lrll}
\toprule
    \textbf{Parcela} & \textbf{Valor} & \textbf{Unidade} & \textbf{Obtenção} \\
\midrule
    Peso específico da MLC        & 4,1202 & kN/m\textsuperscript{3} & $\rho_{mean} = 420$ kg/m\textsuperscript{3} \\
    Volume de uma viga            & 12,2644 & m\textsuperscript{3} & Integração de $A(x)$ \\
    Peso linear da viga           & 2,4954 & kN/m & Calculado \\
    Viga, por área de cobertura   & 0,4340 & kN/m\textsuperscript{2} & Calculado \\
    Terças, por área de cobertura & 0,1889 & kN/m\textsuperscript{2} & Calculado \\
    Telha metálica                & 0,1000 & kN/m\textsuperscript{2} & Declarado \\
    Instalações e serviços        & 0,1000 & kN/m\textsuperscript{2} & Declarado \\
\midrule
    \textbf{Total sobre a viga}   & \textbf{0,8229} & \textbf{kN/m\textsuperscript{2}} & \\
    Total sobre a terça           & 0,3889 & kN/m\textsuperscript{2} & \\
\bottomrule
\end{tabular}}{Elaborado pelo autor (2026)}{peso-proprio}{}{}

A distinção entre parcelas calculadas e declaradas importa. O total de
0,8229 kN/m\textsuperscript{2} é cerca de metade do valor de
1,50 kN/m\textsuperscript{2} que uma estimativa global usual atribuiria a uma
cobertura desse porte, e essa diferença reflete a leveza própria da solução em
madeira, que se perderia caso se adotasse valor único de tabela.

\subsection{Sobrecarga}\label{ssec:sobrecarga}

A cobertura enquadra-se na categoria H da EN 1991-1-1 \cite{EN1991-1-1:2002},
que abrange coberturas não acessíveis salvo para manutenção e reparo. Adotou-se
sobrecarga distribuída de 0,40 kN/m\textsuperscript{2} e carga concentrada de
1,00 kN, esta última aplicada apenas na verificação local das terças e da telha.

O enquadramento na categoria H tem efeito determinante nas combinações: os
coeficientes de combinação $\psi_0$, $\psi_1$ e $\psi_2$ da sobrecarga são
nulos, conforme registrado na \autoref{tab:ndp}. Disso resulta que a sobrecarga
e o vento nunca atuam simultaneamente, e que cada um deles governa combinações
próprias.

\subsection{Ação do vento}\label{ssec:acao-vento}

Conforme exposto na \autoref{ssec:rota-vento}, a cobertura não se enquadra como
cobertura isolada, por ser telhado múltiplo e por haver obstrução permanente sob
ela. Aplicou-se a rota de edificação com coeficientes de pressão externa e
interna da ABNT NBR 6123 \cite{NBR6123:2023}.

Para o vento perpendicular à cumeeira, foram varridos oito padrões de
coeficientes, combinando os dois tramos da cobertura, as duas direções de
incidência e os dois valores admissíveis de coeficiente de pressão interna,
$+0{,}8$ e $-0{,}3$. A varredura é completa, isto é, nenhum padrão foi descartado a
priori por julgamento de qual seria o mais desfavorável.

Para o vento paralelo à cumeeira, a edificação aberta não oferece superfície de
barlavento, e a ação resulta do arrasto nos pilares somado ao atrito sobre a
cobertura. A \autoref{tab:acoes-vento} resume as duas direções, com os valores
obtidos conforme a ABNT NBR 6123 \cite{NBR6123:2023} e a parcela de atrito
calculada pelo item 7.2.2, com $a \times b = 1.863$ m\textsuperscript{2} em
planta.

\tabela{Ações de vento por direção de incidência}{\footnotesize
\begin{tabular}{lrll}
\toprule
    \textbf{Grandeza} & \textbf{Valor} & \textbf{Unidade} & \textbf{Observação} \\
\midrule
    \multicolumn{4}{l}{\textit{Vento perpendicular à cumeeira}} \\
    Padrões de coeficientes varridos & 8 & & Dois tramos, duas direções, dois $c_{pi}$ \\
    Combinações de ELU com vento & 17 & & Varredura completa \\
    Sucção governante & $-7{,}609$ & kN/m & Padrão com $c_{pi} = +0{,}8$ \\
\addlinespace
    \multicolumn{4}{l}{\textit{Vento paralelo à cumeeira}} \\
    Arrasto nos pilares & 60,5 & kN & 27 pilares \\
    Atrito sobre a cobertura & 776,8 & kN & $F_t = 0{,}5\,q\,a\,b$ \\
    \textbf{Força longitudinal total} & \textbf{837,3} & \textbf{kN} & \\
    Força distribuída equivalente & 18,202 & kN/m & \\
\bottomrule
\end{tabular}}{Elaborado pelo autor (2026)}{acoes-vento}{}{}

A parcela de atrito merece comentário, por ser a que governa todo o
dimensionamento do contraventamento longitudinal. Ela equivale a 2,35 vezes o
arrasto que se obteria pelo bloqueio da silhueta da edificação, o que à primeira
vista sugere erro de cálculo. Não é o caso: a prescrição do item 7.2.2 é
envolvente por construção, pois reúne em uma só expressão a ação sobre ambas as
faces da cobertura. A ordem de grandeza é característica da prescrição adotada,
e foi mantida por ser a que a norma determina.

\section{COMBINAÇÕES DE AÇÕES}\label{sec:combinacoes}

As combinações seguem a EN 1990 \cite{EN1990:2002}, com os coeficientes parciais
e de combinação registrados na \autoref{tab:ndp}. A
\autoref{tab:combinacoes} descreve os grupos considerados. Não há termo de
sobrecarga nas combinações com vento, nem termo de vento nas combinações com
sobrecarga, por serem nulos os coeficientes $\psi$ da categoria H. A combinação
de sucção adota o coeficiente favorável de 0,9 para a ação permanente, conforme
a tabela A1.2(A), para verificações de equilíbrio.

\tabela
{Grupos de combinações de ações considerados}
{\footnotesize
\begin{tabular}{llp{5.4cm}}
\toprule
    \textbf{Grupo} & \textbf{Expressão} & \textbf{Finalidade} \\
\midrule
    Sobrecarga dominante & $1{,}35\,G + 1{,}5\,Q$ & Momento e cortante máximos na viga \\
    Vento de pressão     & $1{,}35\,G + 1{,}5\,W$ & Flexão com vento descendente \\
    Vento de sucção      & $0{,}9\,G + 1{,}5\,W$  & Arrancamento, ancoragem e estabilidade \\
    Acidental            & $G + A_d$          & Impacto de veículo em pilar \\
    Execução             & $\gamma_G\,\varphi_{dyn}\,G$ & Içamento da viga na montagem \\
    Serviço              & Característica e quase permanente & Flecha instantânea e final \\
\bottomrule
\end{tabular}}
{Elaborado pelo autor (2026)}
{combinacoes}
{}
{}

Os oito padrões de coeficientes de pressão geram, cada um, uma combinação de
pressão e uma de sucção. Somadas à combinação de sobrecarga dominante, resultam
dezessete combinações de estado-limite último, todas efetivamente avaliadas.

\section{ESFORÇOS SOLICITANTES}\label{sec:esforcos}

A \autoref{tab:esforcos-governantes} apresenta os esforços de cálculo que
governam o dimensionamento da viga, nas combinações da
EN 1990 \cite{EN1990:2002}. O esforço normal reduzido decorre do modelo
biapoiado, discutido na \autoref{ssec:modelo-apoio}.

\tabela{Esforços de cálculo governantes na viga}{\footnotesize
\begin{tabular}{lcrl}
\toprule
    \textbf{Esforço} & \textbf{Símbolo} & \textbf{Valor} & \textbf{Combinação governante} \\
\midrule
    Momento fletor na cumeeira & $M_{Ed}$ & 504,26 kN$\cdot$m & $1{,}35\,G + 1{,}5\,Q$ \\
    Esforço cortante           & $V_{Ed}$ &  99,61 kN         & $1{,}35\,G + 1{,}5\,Q$ \\
    Esforço normal             & $N_{Ed}$ &   8,75 kN         & $1{,}35\,G + 1{,}5\,Q$ \\
    Esforço horizontal na base & $H_{Ed}$ &   0,00 kN         & Modelo biapoiado \\
    Carga de sucção            & $q_{Ed}$ & $-7{,}61$ kN/m      & $0{,}9\,G + 1{,}5\,W$ \\
\bottomrule
\end{tabular}}{Elaborado pelo autor (2026)}{esforcos-governantes}{}{}

A \autoref{fig:modelo-momento} apresenta a distribuição do momento fletor ao
longo das vigas, na combinação governante, em escala de kN$\cdot$m. Terças e
pilares aparecem em cinza, por não ser o momento a grandeza que os dimensiona. O
modelo de elementos finitos reproduz a condição de apoio adotada, sem restrição
horizontal, e devolve momento máximo de 504,2 kN$\cdot$m contra os
504,26 kN$\cdot$m da solução analítica.

\figura
{Momento fletor nas vigas principais na combinação governante de estado-limite último}
{0.98}
{fig/modelo_momento.png}
{Elaborado pelo autor (2026)}
{modelo-momento}
{}
{}

Registre-se um resultado que contraria a expectativa usual para coberturas
leves. A combinação que governa o momento não é uma combinação de vento, e sim a
de sobrecarga dominante. Ainda que a sucção do vento produza carga distribuída
de módulo comparável ao do carregamento gravitacional, ela atua em sentido
oposto ao peso próprio e é parcialmente cancelada por ele. A sucção permanece
determinante, porém, para o arrancamento e para a ancoragem, tratados no
\autoref{cap:detalhamento}.

O esforço normal na viga, de apenas 8,75 kN, é consequência direta do modelo de
apoio adotado. Sob o modelo de arco, esse valor seria de ordem de grandeza
superior, e a verificação de flexo-compressão teria peso bem maior no
dimensionamento.

\section{ORDEM DA ANÁLISE E ESTABILIDADE GLOBAL}\label{sec:ordem-analise}

Resta estabelecer se a análise de primeira ordem é suficiente ou se os efeitos
de segunda ordem devem ser incorporados. O critério é o fator $\alpha_{cr}$,
razão entre a carga crítica elástica e a carga de cálculo, com limite de 10
conforme o item 5.2.1 da EN 1993-1-1 \cite{EN1993-1-1:2005}.

Na direção longitudinal, obteve-se $\alpha_{cr} = 148{,}47$. Na direção
transversal, mais flexível por depender do diafragma de cobertura e dos
frontões, obteve-se $\alpha_{cr} = 20{,}88$. O menor dos dois, 20,88, supera o
limite com folga superior ao dobro, de modo que o sistema é indeslocável nas
duas direções e a análise de primeira ordem é suficiente. O fator de
amplificação calculado, de 1,0503, não foi aplicado.

As imperfeições geométricas foram consideradas por força horizontal equivalente,
com $\phi_0 = 0{,}005$, do que resulta excentricidade inicial de 22,91 mm e
momento de 2,2821 kN$\cdot$m nos pilares, correspondente a taxa de utilização de
0,047. Os efeitos de segunda ordem ao nível de barra, por sua vez, não são
desprezados: permanecem incorporados pelos coeficientes $k_c$ das verificações
de flexo-compressão, conforme o item 6.3.2 da
EN 1995-1-1 \cite{EN1995-1-1:2004}.
